% Section 7.8, from lectures/L07/appendix/b_exercises.md.

\section{Uppgifter}
\label{c7:sec:uppgifter}

\begin{aside}
\textbf{Så arbetar du med uppgifterna.} Vektoruppgifterna i del~1 räknas under lektionen: först
på egen hand, några minuter per uppgift, sedan gemensamt. Del~2 är extrauppgifter att träna
vidare på efter lektionen. De flesta går att räkna i huvudet eller med papper och penna.

\facitnotis{L07}
\end{aside}

% ------------------------------------------------------------------------------------------
\uppgiftsdel{Del 1}{Vektoruppgifter}

\uppgift{1.1}{Addition och skalärmultiplikation}
Vektorerna $\mathbf{u} = (3;\, 2)$ och $\mathbf{v} = (4;\, 6)$ är givna. Bestäm koordinaterna
för:
\begin{delar}(4)
  \task $\mathbf{u} + \mathbf{v}$
  \task $2\mathbf{u} + 3\mathbf{v}$
  \task $\mathbf{u} - \mathbf{v}$
  \task $-2\mathbf{u} - 4\mathbf{v}$
\end{delar}

\uppgift{1.2}{Räkning med tre vektorer}
Vektorerna $\mathbf{a} = (1;\, 4)$, $\mathbf{b} = (-4;\, 2)$ och $\mathbf{c} = (3;\, -4)$ är
givna. Bestäm koordinaterna för:
\begin{delar}(4)
  \task $\mathbf{a} + \mathbf{b} + \mathbf{c}$
  \task $\mathbf{a} - \mathbf{b} - \mathbf{c}$
  \task $2\mathbf{a} + 4\mathbf{b}$
  \task $2\mathbf{a} - 4\mathbf{b}$
\end{delar}

\uppgift{1.3}{Absolutbelopp}
Vektorerna $\mathbf{a} = (2;\, 1)$ och $\mathbf{b} = (1;\, -2)$ är givna. Bestäm i exakt form
längden av:
\begin{delar}(4)
  \task $\mathbf{a}$
  \task $3\mathbf{a} + \mathbf{b}$
  \task $\mathbf{a} - \mathbf{b}$
  \task $3\mathbf{a} - 4\mathbf{b}$
\end{delar}

\uppgift{1.4}{Lös ekvationerna}
Bestäm $x$ och $y$.
\begin{delar}(2)
  \task $(x, y) = 5(3, 1) + 2(4, -5)$
  \task $(15, 12) = (x, 10) + (23, y)$
\end{delar}

\uppgift{1.5}{Parallella vektorer}
Om $\mathbf{u} = k \cdot \mathbf{v}$ är vektorerna $\mathbf{u}$ och $\mathbf{v}$ parallella. Vilka
av följande vektorer är parallella med $\mathbf{u} = (2;\, -3)$?
\begin{delar}(4)
  \task $(6;\, -9)$
  \task $(-8;\, 16)$
  \task $(10;\, 15)$
  \task $(-12;\, 18)$
\end{delar}

\uppgift{1.6}{Bestäm okända koefficienter}
Vektorerna $\mathbf{u} = (-2;\, 3)$ och $\mathbf{v} = (3;\, -1)$ är givna. Bestäm talen $a$ och
$b$ så att
\[
  a\mathbf{u} + \mathbf{v} = (0;\, b).
\]

% ------------------------------------------------------------------------------------------
\uppgiftsdel{Del 2}{Extrauppgifter}
Uppgifterna nedan är extra träning och görs med fördel på egen hand efter lektionen.

\uppgift{2.1}{Avstånd mellan punkter}
Beräkna avståndet mellan punkterna:
\begin{delar}(3)
  \task $(0, 0)$ och $(3, 4)$
  \task $(1, 2)$ och $(4, 6)$
  \task $(-2, 1)$ och $(2, -2)$
  \task $(5, -1)$ och $(5, 7)$
  \task $(-3, 2)$ och $(1, 2)$
\end{delar}

\uppgift{2.2}{Vektorn mellan två punkter}
Vektorn från punkten $A$ till punkten $B$ ges av
$\overrightarrow{AB} = (x_B - x_A;\, y_B - y_A)$.
\begin{parts}
  \item $A = (1, 2)$ och $B = (5, 5)$. Bestäm $\overrightarrow{AB}$ och
    $\abs{\overrightarrow{AB}}$.
  \item $A = (-1, 4)$ och $B = (2, 0)$. Bestäm $\overrightarrow{AB}$ och
    $\abs{\overrightarrow{AB}}$.
  \item Bestäm $\overrightarrow{BA}$ i a) och jämför med $\overrightarrow{AB}$.
  \item Vad gäller allmänt mellan $\overrightarrow{AB}$ och $\overrightarrow{BA}$?
\end{parts}

\uppgift{2.3}{Absolutbelopp}
Bestäm $\abs{\mathbf{u}}$ i exakt form och som närmevärde med två decimaler.
\begin{delar}(3)
  \task $\mathbf{u} = (6;\, 8)$
  \task $\mathbf{u} = (-5;\, 12)$
  \task $\mathbf{u} = (1;\, 1)$
  \task $\mathbf{u} = (-3;\, -3)$
  \task $\mathbf{u} = (0;\, -7)$
\end{delar}

\uppgift{2.4}{Vinkel mot $x$-axeln}
Bestäm vinkeln $v$ mot den positiva $x$-axeln. Ange vilken kvadrant vektorn ligger i och om
kvadrantkorrigering behövs.
\begin{delar}(5)
  \task $(1;\, 1)$
  \task $(-1;\, 1)$
  \task $(-1;\, -1)$
  \task $(1;\, -1)$
  \task $(4;\, 3)$
\end{delar}

\uppgift{2.5}{Från belopp och vinkel till komponenter}
En vektor med längden $\abs{\mathbf{u}}$ och vinkeln $v$ mot den positiva $x$-axeln har
komponenterna
\[
  u_x = \abs{\mathbf{u}} \cos v, \qquad u_y = \abs{\mathbf{u}} \sin v.
\]
Bestäm komponenterna:
\begin{delar}(2)
  \task $\abs{\mathbf{u}} = 10$, $v = 30^{\circ}$
  \task $\abs{\mathbf{u}} = 5$, $v = 90^{\circ}$
  \task $\abs{\mathbf{u}} = 4$, $v = 180^{\circ}$
  \task $\abs{\mathbf{u}} = 2$, $v = 225^{\circ}$
  \task* Kontrollera i a) att $\sqrt{u_x^2 + u_y^2} = 10$.
\end{delar}

\uppgift{2.6}{Skalärprodukt}
Beräkna $\mathbf{u} \cdot \mathbf{v}$.
\begin{delar}(2)
  \task $\mathbf{u} = (2;\, 3)$, $\mathbf{v} = (4;\, 1)$
  \task $\mathbf{u} = (1;\, -2)$, $\mathbf{v} = (3;\, 5)$
  \task $\mathbf{u} = (2;\, 4)$, $\mathbf{v} = (-6;\, 3)$
  \task $\mathbf{u} = (5;\, 0)$, $\mathbf{v} = (0;\, 7)$
  \task* Vad har c) och d) gemensamt, och vad säger det om vektorerna?
\end{delar}

\uppgift{2.7}{Vinkelräta vektorer}
Två vektorer är vinkelräta (ortogonala) om $\mathbf{u} \cdot \mathbf{v} = 0$.
\begin{parts}
  \item Är $(3;\, 4)$ och $(-4;\, 3)$ vinkelräta?
  \item Är $(1;\, 2)$ och $(2;\, 1)$ vinkelräta?
  \item Bestäm talet $k$ så att $(2;\, k)$ blir vinkelrät mot $(6;\, -3)$.
  \item Ange en vektor som är vinkelrät mot $(5;\, -2)$.
\end{parts}

\uppgift{2.8}{Vinkeln mellan två vektorer}
Beräkna vinkeln $\alpha$ mellan vektorerna med hjälp av
$\cos\alpha = \dfrac{\mathbf{u} \cdot \mathbf{v}}{\abs{\mathbf{u}}\abs{\mathbf{v}}}$.
\begin{delar}(2)
  \task $\mathbf{u} = (1;\, 0)$, $\mathbf{v} = (1;\, 1)$
  \task $\mathbf{u} = (3;\, 0)$, $\mathbf{v} = (0;\, 4)$
  \task $\mathbf{u} = (2;\, 1)$, $\mathbf{v} = (-1;\, 3)$
  \task $\mathbf{u} = (2;\, 2)$, $\mathbf{v} = (-3;\, -3)$
\end{delar}

\uppgift{2.9}{Enhetsvektorer}
En enhetsvektor har längden $1$ och fås som
$\hat{\mathbf{u}} = \dfrac{\mathbf{u}}{\abs{\mathbf{u}}}$.
\begin{parts}
  \item Bestäm enhetsvektorn i riktning $(3;\, 4)$.
  \item Kontrollera att svaret i a) verkligen har längden $1$.
  \item Bestäm enhetsvektorn i riktning $(-6;\, 8)$.
  \item Bestäm en vektor med längden $15$ i samma riktning som $(3;\, 4)$.
  \item Bestäm en vektor med längden $15$ som är motriktad $(3;\, 4)$.
\end{parts}

\uppgift{2.10}{Resultant}
Tre krafter verkar i samma punkt: $\mathbf{F}_1 = (4;\, 2)$, $\mathbf{F}_2 = (-1;\, 5)$ och
$\mathbf{F}_3 = (-3;\, -7)$.
\begin{parts}
  \item Beräkna resultanten $\mathbf{F} = \mathbf{F}_1 + \mathbf{F}_2 + \mathbf{F}_3$.
  \item Vad innebär resultatet?
  \item Beräkna resultanten och dess belopp om $\mathbf{F}_3$ tas bort.
  \item Vilken vektor måste adderas till $\mathbf{F}_1 + \mathbf{F}_2$ för att resultanten ska
    bli noll?
\end{parts}

\uppgift{2.11}{Parallella och motriktade vektorer}
Vektorn $\mathbf{u} = (4;\, -6)$ är given.
\begin{parts}
  \item Vilka av $(2;\, -3)$, $(-8;\, 12)$, $(6;\, -4)$ och $(-2;\, 3)$ är parallella med
    $\mathbf{u}$?
  \item Vilka av dem är motriktade $\mathbf{u}$?
  \item Bestäm $t$ så att $(t;\, 9)$ blir parallell med $\mathbf{u}$.
  \item Kan en vektor vara parallell med sig själv? Motivera.
\end{parts}

\uppgift{2.12}{Spänningar som visare}
I ett växelströmsnät representeras spänningar av vektorer, så kallade \textbf{visare}. Över ett
motstånd och en spole i serie ligger spänningarna $\mathbf{U}_R = (12;\, 0)\,\text{V}$ och
$\mathbf{U}_L = (0;\, 5)\,\text{V}$.
\begin{parts}
  \item Beräkna den totala spänningen $\mathbf{U} = \mathbf{U}_R + \mathbf{U}_L$.
  \item Beräkna beloppet $\abs{\mathbf{U}}$.
  \item Beräkna vinkeln mot den positiva $x$-axeln.
  \item Varför är $\abs{\mathbf{U}}$ inte lika med $12 + 5 = 17\,\text{V}$?
\end{parts}

\uppgift{2.13}{Skalärproduktens räkneregler}
Vektorerna $\mathbf{u} = (2;\, -1)$, $\mathbf{v} = (3;\, 4)$ och $\mathbf{w} = (-1;\, 2)$ är
givna.
\begin{parts}
  \item Beräkna $\mathbf{u} \cdot \mathbf{v}$ och $\mathbf{v} \cdot \mathbf{u}$. Vad gäller?
  \item Beräkna $\mathbf{u} \cdot (\mathbf{v} + \mathbf{w})$ och
    $\mathbf{u} \cdot \mathbf{v} + \mathbf{u} \cdot \mathbf{w}$. Vad gäller?
  \item Beräkna $\mathbf{u} \cdot \mathbf{u}$ och jämför med $\abs{\mathbf{u}}^2$.
  \item Vad blir $\mathbf{u} \cdot \mathbf{u}$ uttryckt i $\abs{\mathbf{u}}$ allmänt?
\end{parts}

\uppgift{2.14}{Hitta felet}
Varje rad innehåller ett vanligt vektorfel. Förklara felet och ange det korrekta svaret.
\begin{parts}
  \item $\abs{(3;\, 4)} = 3 + 4 = 7$
  \item $2(3;\, -1) = (6;\, -1)$
  \item $(2;\, 3) \cdot (4;\, 5) = (8;\, 15)$
  \item Vektorn $(-3;\, 4)$ bildar vinkeln
    $\arctan\!\left(\dfrac{4}{-3}\right) \approx -53{,}1^{\circ}$ mot den positiva $x$-axeln.
  \item $\abs{-2\mathbf{u}} = -2\abs{\mathbf{u}}$
\end{parts}
